\documentclass[11pt]{article}

\newcommand{\cnum}{Math 245B}
\newcommand{\ced}{Winter 2026}
\newcommand{\ctitle}[4]{\title{\vspace{-0.5in}\cnum, \ced\\Assignment #1: #2}\author{\vspace{-0.35in}\\#3, #4}}
\usepackage{enumitem}
\usepackage[usenames,dvipsnames,svgnames,table,hyperref]{xcolor}
\usepackage{amsmath, amsfonts}
\usepackage{graphicx} % Required for inserting images
\usepackage{float}
\usepackage{listings}
\usepackage{xcolor}
\usepackage{hyperref}
\usepackage{comment}

\renewcommand*{\theenumi}{\alph{enumi}}
\renewcommand*\labelenumi{(\theenumi)}
\renewcommand*{\theenumii}{\roman{enumii}}
\renewcommand*\labelenumii{\theenumii.}
\newcommand{\notiff}{%
  \mathrel{{\ooalign{\hidewidth$\not\phantom{"}$\hidewidth\cr$\iff$}}}}

\definecolor{codegreen}{rgb}{0,0.6,0}
\definecolor{codegray}{rgb}{0.5,0.5,0.5}
\definecolor{codepurple}{rgb}{0.58,0,0.82}
\definecolor{backcolour}{rgb}{0.95,0.95,0.92}
\definecolor{header}{HTML}{205459}
\definecolor{solution}{HTML}{204d52}

\newcommand{\solution}[1]{{{\textcolor{header}{Solution:} \textcolor{solution}{#1}}}}
\setlist[enumerate]{label=(\alph*)}

\lstdefinestyle{mystyle}{
    backgroundcolor=\color{backcolour},
    commentstyle=\color{codegreen},
    keywordstyle=\color{magenta},
    numberstyle=\tiny\color{codegray},
    stringstyle=\color{codepurple},
    basicstyle=\ttfamily\footnotesize,
    breakatwhitespace=false,
    breaklines=true,
    captionpos=b,
    keepspaces=true,
    numbers=left,
    numbersep=5pt,
    showspaces=false,
    showstringspaces=false,
    showtabs=false,
    tabsize=2
}


\begin{document}
\ctitle{\#1}{}{Asher Christian}{006-150-286}
\date{}
\maketitle

\section{martingale convergence theorem}
Example
\[
    \{X_n\}
\] 
is martingale simple random walk or smething where you buy stocks and you want to know how many times with $X$ go from  $a$ to  $b$
before it goes to 0.
Suppose  $X_n \ge -M$ for any  $n$ almost sure. then
 \[
Y_n \ge (\# a \nearrow b) |b-1| - (a +M)
\] 
and $ \mathbb{E}Y_n = 0$
this implies that
\[
\mathbb{E}(\# a \nearrow b) \le \frac{a+m}{b-a}
\] 
and
\[
    \mathbb{E}(\# b \searrow a) \le  \frac{(M-b)}{b-a} \qquad \text{ if } \sup_{n}|X_n| \le M
\] 
so for martingale $\{X_n\}$ with uniform lower bound for all $a,b$  $ \mathbb{E}(\# a \nearrow b) < \infty \implies \forall a,b$ $ \# a \nearrow b < \infty$ almost sure.
\[
   \mathbb{P}( \bigcap_{a,b \in \mathbb{Q}}\{ \# a \nearrow b < \infty\} ) = 1
\] 
if $\limsup_{n\to \infty}X_n > \liminf_{n\to \infty}X_n$ then there exists $a,b \in \mathbb{Q}$ such that $\# a \ne b = +\infty$.
If for some $\Omega$
 \[
\limsup_{n\to \infty} X_n(\omega) > \liminf_{n\to \infty} X_n(\omega)
\] 
then for this omega there exists $a,b \in \mathbb{Q}$ such that
\[
    \# a \ne b (\{X_n(\omega)\}) = +\infty
\] 
so we have
\[
\mathbb{P}(\limsup_{n\to \infty} X_n > \liminf_{n\to \infty}X_n) = 0
\] 
then we need to know
\[
\mathbb{P}(\limsup_{n\to \infty}X_n = +\infty)
\] 
but
\[
\lim_{k\to \infty} \mathbb{P}(X_n \ge k) = 0
\] 
\[
    \mathbb{E}[X_n] = 0 \implies \mathbb{E}[X_n^{+}] = \mathbb{E}[X_n^{-}] \qquad \mathbb{E}[X_n^{-}] \le M
\] 
putting together we get the $\lim_{n\to \infty}X_n$ exists almost sure.
General case 
\section{MCT}
\emph{
    $\{X_n\}$ is a sub martingale and  $\sup_{n} \mathbb{E}X_n^{+} < \infty$  then
    \[
        \lim_{n\to \infty}X_n \rightarrow X_\infty \qquad \text{a.s}
    \] 
}
for simplicity choose $a = 0, b =1$  $\{X_n\}$ is sub martingale then  $X_n^{+}$ is sub martingale
since if $\phi$ is convex and non decreasing then $\phi(X_n)$ is a sub martingale, in this case $\phi(x) = \max(0,x)$.
And
\[
Y_n = (A \cdot X^{+})_n \qquad Z_n = (B \cdot X^{+})_n  \qquad B_n = 1 - A_n
\] 
then $Y_n$ and  $Z_n$ are also sub martingales since  $A_n, B_n \ge 0$ 
\[
    \mathbb{E}[Y_n] \ge \mathbb{E}Y_0 = 0, \qquad \mathbb{E}Z_n \ge \mathbb{E}Z_0 = 0
\] 
and $Y_n + Z_n = X_n^{+} - X_0^{+}$ so
\[
    \mathbb{E}[Y_n] + \mathbb{E}[Z_n] = \mathbb{E}[X_n^{+}] - \mathbb{E}[X_0] \le M
\] 
so
\[
    \mathbb{E}[Y_n] \le \mathbb{E}[X_{n}^{+}] - \mathbb{E}[X_0^{+}] \le M
\] 
and $Y_n \ge \#a \nearrow b (b-a)$ so
 \[
     (\# a \nearrow b) < \infty \qquad \text{ a.s.}
\] 
Example: POlya's Urn
\[
    \{X_n\}
\] 
is a martingale and $|X_n| \le 1$ a.s. Start with  $R_0 = B_0 = 1$ and $X_n = \frac{R_n}{R_n + B_n}$ 
so with martingale convergence theorem there exists $X_\infty$ such that $X_n \rightarrow X_\infty$ almost sure
we know almost sure convergence implies convergence in distribution so $X_n \rightarrow^{d} X_\infty$
It is easy to check
\[
\mathbb{P}(X_n = \frac{m}{n+2}) = \frac{1}{n+1} \qquad \forall 1 \le m \le n + 1, m \in  \mathbb{Z}
\] 
and so $X_\infty \sim^{d} unif[0,1]$.
How about $L_p$ convergence, i.e.  $ \mathbb{E}|X_n - X_\infty|^{p} \rightarrow 0?$ we know
$a.s.$ implies in prob implies in distributions and  $L_p$ implies in probability.
For this question,  $|X_n| \le 1 \implies |X_n - X_\infty| \le 1$ so with BCT we have
\[
\mathbb{E}|X_n - X_\infty|^{p} \rightarrow 0
\] 
Examples $\{X_n\}$ martingale  $\sup_n |X_{n+1} - X_n| \le M$ for some  $M$ 
\[
    A = \{\lim_{n\to  \infty}X_n \in (-\infty,\infty)\} = \{ \limsup_{n\to \infty}X_n = \liminf_{n\to \infty}X_n \in (-\infty,\infty)\}
\] 
\[
    B = \{\limsup_{n\to \infty}X_n = +\infty , \liminf_{ n\to \infty}X_n = -\infty\}
\] 
then
\[
\mathbb{P}(A \cup B) = 1
\] 
in the srw case, $ \mathbb{P}(B) = 1$ Let $K$ be a large number, define
 \[
     T = \min\{n : X_n \le -k\}
\] 
and let
\[
    Y_n = X_{n \wedge T}
\] 
then
\[
Y_n \ge - - M
\] 
so with MCT, $Y_n \rightarrow Y_\infty$ almost sure this implies that $ A \supset \{\inf X_n \ge -k\}$ it holds for any $k \ge 0$ 
 \[
     \mathbb{P}(A \cup \{\liminf_{ X_n\to  -\infty\}} = 1) \implies \mathbb{P}(\{\liminf_{n\to \infty} X_n = -\infty\} \triangle A^{c}) = 0
\] 
This construction is symmetric so we can say that
\[
    \mathbb{P}(\{\limsup_{n\to \infty}X_n = +\infty\} \triangle A^{c}) = 0
\] 
so
\[
    \mathbb{P}(A \cup B) = 1
\] 
In $L_1$ simple random walk, $X_0 = 10$ and  $T = \min\{n: X_n = 0\}$ and  $Y_n = X_{n \wedge T}, Y_n \ge 0, Y_n \rightarrow Y_\infty$ and $Y_\infty = 0$ but $ \mathbb{E}|Y_n - Y_\infty| = \mathbb{E}[Y_n] = 10$
So there is no $L_1$ convergence.
\section{ Martingale $L_p$ Convergence theorem}
Question: $ \mathbb{E}|X_n - X_\infty|^{p} \rightarrow 0$ \\
Theorem: $\{X_n\}$ is martingale and $\sup_n \mathbb{E} |X_n|^{p}< \infty$ then
\[
    \exists X_ \infty \qquad    X_n \rightarrow X_{\infty} \text{ in $L_p$ } p > 1
\] 

\section{Friday}
\emph{
    let $\{X_n\}$ martingale and  $\sup_n \mathbb{E}|X_n|^{p} \le M < \infty$ for $p > 1$ then
    there exists  $X_{\infty} \in \mathcal{F}_\infty = \sigma(\bigcup \mathcal{F}_n)$ such that $X_n \rightarrow X_{\infty}$ in $L_p$
    and  $ \mathbb{E}|X_n - X_\infty|^{p } \rightarrow 0$
}
Proof: First we have $\sup_n \mathbb{E}|X_n| < \infty$ since $p > 1$
together with  $\{X_n\}$ martingale implies that there exists  $X_\infty$ such that $X_n \rightarrow X_\infty$ almost sure by the martingale convergence theorem
\[
    Y_n = X_n - X_\infty \qquad Y_n \rightarrow 0 \text{ a.s.}
\] 
we hope to see that $ \mathbb{E}\sup_n |X_n - X_\infty|^{p} < \infty$ if so we can use dominated convergence to imply that $X_n - X_\infty \rightarrow 0$ in $L_p$
Take a look at
 \[
|X_n - X_\infty|^{p} \le  |X_n|^{p} + |X_\infty|^{p}
\] 
so
\[
\sup_n |X_n - X_\infty|^{p} \le (\sup_n |X_n|)^{p} + |X_\infty|^{p}
\] 
and
\[
\mathbb{E} \sup_n |X_n- X|^{p} \le \mathbb{E}(\sup_n |X_n|)^{p} + \mathbb{E}|X_\infty|^{p}
\] 
\[
\mathbb{E} \sup_n |X_n|^{p} < \infty ?
\] 
\section{theorem Doob's maximum Inequality}
if $\{X_n\}$ is a nonnegative sub martingale, $p > 1$ then for  $X_n^{*} := \max_{1 \le k \le n} X_k$ 
And if $ \mathbb{E}|X_k|^{p} < \infty$ for all $k$
\[
\mathbb{E}|X_n^{*}|^{p} \le \frac{p}{1-p} \mathbb{E}|X_n|^{p}
\] 
\end{document}
